\chapter{Экстремумы функций нескольких переменных. Необходимые условия, достаточные условия.}

\section{Необходимое условие}

%----------------------------------------------------------------------------------------
% Определение 18.5 --- Стационарная точка
%----------------------------------------------------------------------------------------

\begin{defn} \label{df18:5}
Если функция $n$ переменных $f$ дифференцируема в точке $x^0$ и $df(x^0) = 0$, то точка $x^0$ называется стационарной точкой функции $f$.
\end{defn}

\begin{notion}
    Условие стационарности для дифференцируемой функции равносильно одновременному выполнению условий
$$
\frac{\partial f}{\partial x_1}(x^0) = 0, \ldots, \frac{\partial f}{\partial x_n}(x^0) = 0.
$$
Это система $n$ уравнений с $n$ неизвестными для нахождения координат точки $x^0$.
\end{notion}


%----------------------------------------------------------------------------------------
% Теорема 18.5 --- Необходимое условие точки локального экстремума
%----------------------------------------------------------------------------------------

\begin{thm}[необходимое условие точки локального экстремума] \label{th18:5}
Если функция $f$ дифференцируема в точке локального экстремума $x^0$, то $x^0$ "--- стационарная точка $f$.
\end{thm}

\begin{proof}
Зафиксируем $i = 1, \ldots, n$. Если $x^0 = (x_1^0, \ldots, x_n^0)$ "--- точка локального экстремума функции $f(x_1, \ldots, x_n)$, то $x_i^0$ "--- точка локального экстремума того же характера для функции одной переменной $\varphi(x_i) = f(x_1^0, \ldots, x_{i-1}^0, x_i, x_{i+1}^0, \ldots, x_n^0)$, где все переменные, кроме $x_i$, зафиксированы. Тогда $\varphi'(x_i^0) = 0$, т.е. $\frac{\partial f}{\partial x_i}(x_1^0, \ldots, x_i^0, \ldots, x_n^0) = 0$. Это верно для всех $i = 1, \ldots, n$, поэтому $x^0$ "--- стационарная точка.
\end{proof}

Теорема ~\ref{th18:5} является естественным обобщением на многомерный случай одномерной теоремы Ферма ~\ref{ch4n1} (и доказывается непосредственным применением этой теоремы).

\begin{notion}
Условие стационарности не является достаточным условием точки локального экстремума для дифференцируемой функции. Для функции $f(x, y) = x^2 - y^2$ (график которой является гиперболическим параболоидом) точка $(0, 0)$ является стационарной ($f'_x = f'_y = 0$), но не является точкой локального экстремума (в ней график имеет характерный вид седла).
\end{notion} 


\section{Достаточное условие}
\subsection{Квадратичные формы}
%----------------------------------------------------------------------------------------
% Квадратичные формы (введение)
%----------------------------------------------------------------------------------------

Напомним, что квадратичной формой от вектора $x = (x_1, \ldots, x_n)^T$ называется числовая функция
$$
K(x) = K(x_1, \ldots, x_n) = \sum_{i=1}^n b_{ii}x_i^2 + 2\sum_{\substack{i,j=1 \\ i<j}}^n b_{ij}x_ix_j,
$$
где $(b_{ij})_{i,j=1}^n$ "--- симметричная квадратная матрица ($b_{ij} = b_{ji}$, $i, j = 1, \ldots, n$, $b_{ij} \in \bbR$). Если $\alpha \in \bbR$, то $K(\alpha x) = \alpha^2 K(x)$. Квадратичная форма называется:
\begin{enumerate}
    \item[1)] положительно определённой, если $\forall\, x \ne 0 \longrightarrow K(x) > 0$;
    \item[2)] отрицательно определённой, если $\forall\, x \ne 0 \longrightarrow K(x) < 0$;
    \item[3)] неопределённой, если $\exists\, x', x'' \colon K(x') > 0, K(x'') < 0$;
    \item[4)] положительно полуопределённой, если $\forall\, x \longrightarrow K(x) \ge 0$, но $\exists\, x \ne 0 \colon K(x) = 0$;
    \item[5)] отрицательно полуопределённой, если $\forall\, x \longrightarrow K(x) \le 0$, но $\exists\, x \ne 0 \colon K(x) = 0$.
\end{enumerate}
Тождественно нулевая квадратичная форма является одновременно положительно и отрицательно полуопределённой; каждая из остальных квадратичных форм принадлежит ровно одному из указанных пяти классов.

%----------------------------------------------------------------------------------------
% Лемма 18.4 --- Ограниченность квадратичной формы
%----------------------------------------------------------------------------------------

\begin{lemm} \label{lm18:4}
Если квадратичная форма $K(x)$ положительно определена, то $\exists\, C > 0 \colon \forall\, x \longrightarrow K(x) \ge C|x|^2$, где $|x| = \sqrt{x_1^2 + \ldots + x_n^2}$. Если квадратичная форма $K(x)$ отрицательно определена, то $\exists\, C > 0 \colon \forall\, x \longrightarrow K(x) \le -C|x|^2$.
\end{lemm}

Пусть функция $f$ дважды непрерывно дифференцируема в точке $x^0 \in \bbR^n$. Тогда её второй дифференциал в точке $x^0$
$$
d^2f(x^0) = \sum_{i=1}^n \frac{\partial^2 f}{\partial x_i^2} dx_i^2 + 2\sum_{\substack{i,j=1 \\ i<j}}^n \frac{\partial^2 f}{\partial x_i \partial x_j} dx_i dx_j,
$$
где частные производные второго порядка берутся в точке $x^0$. Это квадратичная форма от вектора $\overrightarrow{dx} = (dx_1, \ldots, dx_n)^T$; напомним, что для дважды непрерывно дифференцируемой функции $\frac{\partial^2 f}{\partial x_i \partial x_j} = \frac{\partial^2 f}{\partial x_j \partial x_i}$.

\subsection{Доказательство теоремы}
%----------------------------------------------------------------------------------------
% Теорема 18.6 --- Достаточные условия точки локального экстремума
%----------------------------------------------------------------------------------------

\begin{thm}[достаточные условия точки локального экстремума] \label{th18:6}
Пусть функция $f$ дважды непрерывно дифференцируема в некоторой окрестности стационарной точки $x^0$. Тогда:
\begin{enumerate}
    \item[1)] если $d^2f(x^0)$ "--- положительно определённая квадратичная форма, то $x^0$ "--- точка строгого локального минимума;
    \item[2)] если $d^2f(x^0)$ "--- отрицательно определённая квадратичная форма, то $x^0$ "--- точка строгого локального максимума;
    \item[3)] если $d^2f(x^0)$ "--- неопределённая квадратичная форма, то $x^0$ не является точкой локального экстремума.
\end{enumerate}
\end{thm}

\begin{proof}
1) Применим формулу Тейлора для функций $n$ переменных с остаточным членом в форме Пеано (теорема 10.10). Так как $f \in C^2(U_\delta(x^0))$, то
$$
f(x) = f(x^0) + df(x^0) + \frac{1}{2}d^2f(x^0) + o(\rho^2),
$$
$$
(dx_1, \ldots, dx_n) \to (0, \ldots, 0),
$$
$$
\rho = \sqrt{dx_1^2 + \ldots + dx_n^2}, \quad dx_i = \Delta x_i = x_i - x_i^0, \quad i = 1, \ldots, n,
$$
$$
\overrightarrow{dx} = (dx_1, \ldots, dx_n)^T, \quad \rho^2 = |\overrightarrow{dx}|^2,
$$
$$
o(\rho^2) = \varepsilon(dx_1, \ldots, dx_n)|\overrightarrow{dx}|^2,
$$
где $\lim\limits_{\substack{dx_1 \to 0 \\ \dots \\ dx_n \to 0}} \varepsilon(dx_1, \ldots, dx_n) = 0$. В стационарной точке $df(x^0) = 0$, поэтому
$$
\Delta f(x^0) = f(x) - f(x^0) = \frac{1}{2}d^2f(x^0) + \varepsilon(dx_1, \ldots, dx_n)|\overrightarrow{dx}|^2.
$$

Если $d^2f(x^0)$ "--- положительно определённая квадратичная форма от вектора $\overrightarrow{dx}$, то по лемме ~\ref{lm18:4} 

$d^2f(x^0) \ge C|\overrightarrow{dx}|^2$. Тогда $\Delta f(x^0) \ge \left(\frac{C}{2} + \varepsilon(dx_1, \ldots, dx_n)\right)|\overrightarrow{dx}|^2$. Так как $\lim \varepsilon(dx_1, \ldots, dx_n) = 0$ при $dx_1 \to 0, \ldots, dx_n \to 0$, то
$$
\exists\, \delta > 0 \colon \forall\, \overrightarrow{dx}, \quad 0 < |\overrightarrow{dx}| < \delta \longrightarrow \frac{C}{2} + \varepsilon(dx_1, \ldots, dx_n) > 0.
$$
Значит, $\forall\, x \in \mathring{U}_\delta(x^0)$ выполняется неравенство $\Delta f(x^0) > 0$, т.е. $f(x) > f(x^0)$. Поэтому $x^0$ "--- точка строгого локального минимума.

2) Доказывается аналогично.

3) Пусть теперь $K(\overrightarrow{dx}) = d^2f(x^0)$ "--- неопределённая квадратичная форма. Тогда существует вектор $\vec{z} \ne \vec{0}$ такой, что $K(\vec{z}) > 0$. Рассмотрим всевозможные векторы вида $\overrightarrow{dx} = \lambda\vec{z}$, $\lambda > 0$ (приращения, сонаправленные с $\vec{z}$). Получим
$$
d^2f(x^0) = K(\overrightarrow{dx}) = K(\lambda\vec{z}) = \lambda^2K(\vec{z}) = \lambda^2\beta|\vec{z}|^2,
$$
где $\beta = \frac{K(\vec{z})}{|\vec{z}|^2}$ "--- фиксированное положительное число, так как $\vec{z}$ "--- фиксированный вектор. В соответствующей точке $x = x^0 + \lambda\vec{z}$ (вектор $\lambda\vec{z}$ отложен из точки $x^0$):
\begin{multline*}
\Delta f(x^0) = f(x) - f(x^0) = \frac{1}{2}d^2f(x^0) + \varepsilon(dx_1, \ldots, dx_n) \cdot |\overrightarrow{dx}|^2 = \\
= \frac{\lambda^2\beta|\vec{z}|^2}{2} + \varepsilon(\lambda z_1, \ldots, \lambda z_n) \cdot |\lambda\vec{z}|^2 = \lambda^2|\vec{z}|^2\left(\frac{\beta}{2} + \varepsilon(\lambda z_1, \ldots, \lambda z_n)\right).
\end{multline*}
Так как $\lim \varepsilon(dx_1, \ldots, dx_n) = 0$ при $dx_1 \to 0, \ldots, dx_n \to 0$, то
$$
\exists\, \delta > 0 \colon \forall\, \overrightarrow{dx}, \quad 0 < |\overrightarrow{dx}| < \delta \longrightarrow \frac{\beta}{2} + \varepsilon(dx_1, \ldots, dx_n) > 0.
$$
Векторы вида $\overrightarrow{dx} = \lambda\vec{z}$ удовлетворяют условию $0 < |\overrightarrow{dx}| < \delta$ при $0 < \lambda < \frac{\delta}{|\vec{z}|}$; для них $\Delta f(x^0) > 0$. Значит, найдутся сколь угодно малые по модулю векторы $\overrightarrow{dx}$ такие, что $f(x) > f(x^0)$.

Аналогично, из существования вектора $\vec{z}^* \ne \vec{0}$ такого, что $K(\vec{z}^*) < 0$, следует, что существуют сколь угодно малые по модулю векторы $\overrightarrow{dx}$ такие, что $f(x) < f(x^0)$. Точка $x^0$ не является точкой локального экстремума.
\end{proof}

Заметим, что если $d^2f(x^0)$ "--- полуопределённая квадратичная форма, то теорема ~\ref{th18:6} не даёт ответа на вопрос о наличии и характере экстремума в стационарной точке. Нужно дополнительное исследование, как правило, сводящееся к непосредственному рассмотрению приращения функции в точке.